\section{What, Historically, Is “the Septuagint”?}

The familiar story is compact: a king wanted the Jewish Scriptures for the Library of Alexandria, seventy or seventy-two scholars translated them, and the resulting book became the Bible of the early Church. Every clause preserves part of an ancient memory. As a continuous account, however, it joins events separated by centuries and gives one name to several different objects.

\historythesis{The Septuagint was not one translation of one bound Hebrew Bible by one committee. It began with a Jewish Greek Torah in Ptolemaic Egypt during the third century \textsc{bc}; other books were translated or composed in Greek through later centuries by agents with different aims and source texts. Jewish readers authorized, revised, and debated these writings. Christians inherited them, made them their principal Greek Old Testament, and transmitted them in new book forms and textual combinations. The “Septuagint” is therefore both an originating translation and the history of a growing, pluriform corpus.}
